\pagina{Metodología experimental y reproducibilidad}
\subsection{Etapas y control de versiones}
Se separan campaña inicial, evaluación distribuida y exploraciones cerradas (sección~\ref{sec:exploration}). M6 incluyó comprobación, ajuste, evaluación, escala y verificación: 12, 120, 720, 180 y 12 procesos, respectivamente; 1\,044 en total. Su análisis principal utiliza los 900 primarios de evaluación y escala, sin ampliar el denominador con verificaciones.

La evaluación inicial usó hasta 10\,000 evaluaciones o 30 segundos por solver en sus fases principales, con 30 semillas por grafo y familia. Evaluación cubrió graph48, graph175, graph74 y graph479; escala añadió graph1001. Las instancias de comprobación y ajuste quedaron diferenciadas para evitar atribuir sus mediciones a evaluación final.

El plan distribuido congeló las mejoras descritas con diez semillas, 1001--1010, y hasta 60 segundos por intento, sin límite de evaluaciones ni iteraciones. Termina antes al certificar un ciclo; el límite externo de 120 segundos incluye arranque y carga y controla bloqueos. El plazo cooperativo puede tener una pequeña sobrepasada de instrumentación.

Las 1001 instancias se asignaron sin solapamientos a cinco integrantes: 220 a Arturo, 196 a Zapata y 195 a cada uno de los demás. Cada grafo completo representa seis familias por diez semillas, es decir, 60 combinaciones. La cola completa era una reserva de 60\,060 intentos; no equivale al número efectivamente ejecutado.

\subsection{Entorno y comparabilidad}
Cada integrante conservó el mismo equipo. Los hilos nativos se fijaron a uno por proceso. La evaluación distribuida contiene 7\,563 intentos secuenciales y 1\,243 intentos bajo un coordinador de hasta cuatro procesos en núcleos distintos. El plan fijaba ejecución serial: la concurrencia se registra como desviación de recursos, con código de solver y presupuesto individual conservados. Los perfiles registran CPython 3.14.10 estándar con GIL, arquitectura de 64 bits y NumPy 2.5.4. El entorno y la revisión del código se identificaron antes de los lotes; se preservaron pruebas de preparación y hashes de los componentes.

\begin{table}[H]\centering\small
\caption{Equipos de la evaluación distribuida. CPU y RAM proceden de la preparación declarada y conservada de cada entrega.}
\begin{tabularx}{\linewidth}{lYrY}\toprule
Equipo & Procesador & RAM (GiB) & Sistema \\ \midrule
Arturo & Intel Core i5-12450H & 31.01 & CachyOS, Linux 7.2.9\\
Zapata & Intel Core i5-1334U & 15.30 & CachyOS, Linux 7.3.1\\
César & AMD Ryzen 7 7730U & 14.96 & CachyOS, Linux 7.3.3\\
Eder & AMD Ryzen 5 4600G & 21.31 & Debian 13, Linux 6.12.113\\
Elías & AMD Ryzen 5 7535HS & 15.21 & Windows 11, compilación 26300\\
\bottomrule\end{tabularx}
\end{table}

Las seis familias comparten hardware por grafo; entre equipos cambian las instancias, confundiendo el efecto de máquina y subconjunto. No se normalizan tiempos por frecuencia, núcleos o RAM. Las tablas separan equipo, algoritmo y modo por intento; las medias conjuntas no acreditan aceleración. La campaña histórica utilizó la laptop de Arturo, afinidad a la CPU lógica 2 y un trabajador.

\subsection{Admisión, registros y auditoría}
Cada intento conserva intención de ejecución, parámetros, semilla, stdout y stderr, estado de salida, mejor candidato y auditoría. El coordinador guarda progreso por intento y no reintenta automáticamente interrupciones. La recuperación distingue un primario interrumpido de uno terminado; no convierte una ruta parcial de stdout en éxito.

Para consolidar se comprobaron SHA-256, CRC e inventarios de cinco paquetes; se extrajeron en directorios separados y se ejecutó el agregador oficial. Este volvió a comprobar identidades, presupuesto, snapshot de fuentes, hashes de artefactos, costo y certificado. Las celdas recalculadas de cada integrante coincidieron exactamente con sus resúmenes originales. La agregación no contiene miembros ausentes ni combinaciones duplicadas.

En la evaluación final hay 8\,804 duraciones de proceso disponibles; dos interrupciones no tienen ese dato. Se conservan sin imputación. La suma de 510\,361.334 segundos entre equipos corresponde a 141.767 horas de procesos, no a duración calendario. Los huecos de los historiales de sesión de Elías y Zapata impiden sustituir esa medida por una suma de sesiones supuestamente completa.

\subsection{Cierre parcial y criterios de análisis}
La meta inicial del plan compartido era alcanzar 60 grafos distintos certificados. El responsable indicó una expectativa aproximada de 20 para la entrega; el consolidado certifica 15 y no alcanza ninguna de esas referencias. Se conserva la procedencia de la expectativa de 20 al cierre y no se presenta como objetivo preregistrado.

El corte contiene 149 grafos iniciados, 144 con 60 combinaciones conciliadas y cinco parciales. Para evitar comparar familias con distinto número de semillas en un mismo resumen, la tabla principal entre algoritmos utiliza solo los 144 grafos completos: 1\,440 intentos por familia. Las incidencias de esos grafos siguen en el denominador. Los totales generales incorporan también los 166 intentos de los cinco grafos parciales.

Las 51\,254 combinaciones no ejecutadas quedan fuera del alcance aceptado. Se excluyó graph57--GA--1001 de Elías por no pertenecer a su paquete de 600 resultados auditados.
